\documentclass{article}
\usepackage[UTF8]{ctex}
\usepackage{amsmath}
\usepackage{amssymb}
\usepackage{graphicx}
\usepackage{geometry}
\geometry{a4paper, margin=2cm}

\newcommand{\problemname}[1]{\begin{center}\Large\bfseries #1\end{center}}
\usepackage{enumitem}

\begin{document}
\problemname{数学考试}

一次数学考试后，蒜头君不敢看自己的分数，他把试卷偷偷塞给你，让你判断下他是否及格。

\subsection*{输入格式}
一个整数 $n\text{ }(1\len\le100)$ 表示分数。

\subsection*{输出格式}
\begin{itemize}
  \item 若分数大于等于 $60$ 分，则及格，告诉他通过了，输出\texttt{pass}。

  \item 否则，你不需要输出任何东西，给他一个眼神就行了。

\end{itemize}
\subsection*{样例输入}
\begin{verbatim}
100
\end{verbatim}

\subsection*{样例输出}
\begin{verbatim}
pass
\end{verbatim}


\end{document}
